% Generated draft: requires mathematical and source review.
\documentclass[11pt]{article}
\usepackage[margin=25mm]{geometry}
\usepackage{fontspec}
\setmainfont{Latin Modern Roman}
\usepackage{amsmath,amssymb,mathtools,unicode-math}
\setmathfont{Latin Modern Math}
\usepackage{enumitem,needspace,xurl,hyperref}
\setlength{\parindent}{0pt}
\setlength{\parskip}{5pt}
\setlength{\emergencystretch}{3em}
\newcommand{\N}{\mathbb{N}}
\newcommand{\Z}{\mathbb{Z}}
\newcommand{\Q}{\mathbb{Q}}
\newcommand{\R}{\mathbb{R}}
\newcommand{\C}{\mathbb{C}}
\newcommand{\F}{\mathbb{F}}
\newcommand{\A}{\mathbb{A}}
\newcommand{\PP}{\mathbb{P}}
\newcommand{\E}{\mathbb{E}}
\newcommand{\eps}{\varepsilon}
\newcommand{\dd}{\,\mathrm{d}}
\begin{document}
\section*{UnsolvedMath ID: 30003817}
\textbf{Source problem code:} catalogue entry
\subsection{Definitions and mathematical statement}
The supplied source poses a conceptual research question rather than a fully quantified conjecture. The following framework describes the enumeration data needed to formulate it precisely; the source record does not specify these data.

Let \(\mathcal U=\bigsqcup_{m\geq 0}\mathcal U_m\) be a graded universe of finite structures, with each \(\mathcal U_m\) finite. Suppose that an enumerative interpretation assigns to every \(p\in\mathcal U\) a sequence \((a_n(p))_{n\geq 0}\) of nonnegative integers. In a pattern-avoidance setting, \(a_n(p)\) would count structures of size \(n\) avoiding \(p\), but neither an avoidance relation nor an ambient counting class is supplied here. Relative to the chosen enumerative interpretation, define Wilf equivalence by
\[
 p\sim_{\mathrm W}q
 \quad\Longleftrightarrow\quad
 a_n(p)=a_n(q)\quad\text{for every }n\geq 0.
\]
Write
\[
 U_m=|\mathcal U_m|,
 \qquad
 K_m=\bigl|\mathcal U_m/\!\sim_{\mathrm W}\bigr|
\]
for the number of structures and the number of Wilf-equivalence classes represented at size \(m\), respectively.

\textbf{Research question, subject to unresolved qualifications.} In natural universes of finite structures equipped with the relevant enumeration data, determine how prevalent Wilf equivalences are and how surprising their occurrence should be. In particular, determine whether exponential Wilf collapse is typical, and identify structural mechanisms explaining the prevalence of Wilf-equivalent objects.

The record does not specify the universes, the counting sequences, the precise asymptotic relationship between \(U_m\) and \(K_m\) meant by ``exponential Wilf collapse,'' or a meaning of ``natural,'' ``typical,'' or ``surprising.'' Consequently, this framework does not select a particular exponential bound, probability model, or universal quantifier, and the supplied question cannot yet be expressed as a single precise conjecture.
\subsection{Short English statement}
The question asks why distinct finite structures so often have identical associated counting sequences, and whether an exponential reduction in the number of distinct counting sequences is typical. The intended families, counting rules, meaning of exponential collapse, and notion of typicality remain unspecified. This is a statement-only draft requiring human mathematical review.
\subsection{Sources}
\begin{itemize}
\item Supplied local UnsolvedMath registry record, permanent ID \(30003817\), titled ``Prevalence of Wilf Equivalence and Exponential Collapse''; exact source problem code \texttt{catalogue entry}. The collection is ``Oberwolfach Reports Open Problems.'' No individual proposer is supplied; the proposed year is recorded as 2018.
\item Supplied report citation: \emph{Enumerative Combinatorics (2019)}, Oberwolfach Reports, DOI \texttt{10.4171/owr/2018/23}, pp.~1381--1464. Source URL and supplied literature reference labeled ``original OWR report'': \url{https://doi.org/10.4171/owr/2018/23}. The background describes Oberwolfach Reports as documenting mathematical workshops and their open-problem discussions.
\item Literature assessment dated exactly \texttt{2026-08-21}: the supplied curation labels the problem open, describes it as a broad conceptual question rather than a single formal assertion, and states that the report gives no general theorem answering it. This is supplied dated curation, not an independently verified literature assessment; time-sensitive status and cited claims require independent verification.
\item The registry labels the statement \texttt{corrected\_verified} and explains that notation, typography, and sentence structure were repaired while mathematical content was preserved. This supplied editorial label does not establish independent mathematical verification of the present draft.
\end{itemize}
\end{document}
